\section{A smooth bounded output: final tanh}
\label{sec:smooth-output}
The bounded output need not have corners. The same polynomial bound holds with a final tanh in both likelihood and posterior prediction, provided the all-width comparison class uses that same output convention.

Let $\mathcal T_{C,h}$ consist of functions $f=\tanh(g)$, where $g$ is an arbitrary-coefficient width-$C$, hidden-depth-$h$ ReLU network with a scalar affine readout. Use an unrestricted Gaussian ReLU learner with the same final tanh. All preceding posterior-certificate arguments remain valid: tanh is odd, bounded by one, and 1-Lipschitz. These three properties give, respectively, the prior mean-zero bound, the bounded-loss estimates, and the same covering and Gaussian-perturbation estimates. Its monotonicity also preserves the parameter-count bound: a threshold on $\tanh(z_w(x))$ is a threshold on $z_w(x)$ with the threshold value transformed by $\operatorname{arctanh}$. Squared-loss superlevel sets are unions of at most two such threshold sets.

\begin{theorem}[The same width polynomial with a final tanh]
\label{thm:smooth-output-oracle}
Assume $|Y|\le1$, $dP_X/dU_Q\le\kappa$, and $H\ge\max\{h,H_d\}$. Use a final tanh in both Gaussian likelihood and prediction, with the same independent Gaussian priors and two-certificate calibration. Simultaneously for all $\gamma\ge1$ and $N\ge C+1$, with probability at least $1-\delta$,
\[
 R(\widehat f_n)\le\inf_{f\in\mathcal T_{C,h}}R(f)+\varepsilon
\]
under the same sample polynomial \eqref{eq:all-width-regular-sample}, with additional logarithmic factors only.
\end{theorem}
\begin{proof}
A function $f\in\mathcal T_{C,h}$ is monotone on each affine interval of its raw network along a coordinate line. There are at most $R_0=(C+1)^h$ such intervals, so its total variation on that line is at most $2R_0$.

Fix $0<\eta<1/2$ and define
\[
 F(x)=\operatorname{arctanh}((1-\eta)f(x)),\qquad
 A_\eta=\operatorname{arctanh}(1-\eta)
       \le\tfrac12\log(2/\eta).
\]
This transformed target is bounded by $A_\eta$. It is still monotone on each of the same intervals and hence has coordinate-line total variation at most $2A_\eta R_0$. The interpolation proof only used these amplitude and variation facts. Applying the hierarchical construction to $F/A_\eta$ and multiplying its readout by $A_\eta$ gives a raw ReLU network $q$ with
\[
 \operatorname{width}(q)=O_d(\Delta^{-d}),\quad
 P(q)=O_d(A_\eta R_0\Delta^{-d}),\quad
 \int_Q|q-F|\,dU_Q\le2A_\eta dR_0\Delta.
\]
The resulting bounded output satisfies
\[
 |\tanh(q(x))-f(x)|\le|q(x)-F(x)|+\eta.
\]
Take $\eta=\varepsilon/32$ and a dyadic $\Delta$ within a factor two below $\varepsilon/[128\kappa A_\eta dR_0]$. Density domination gives a population squared-loss increase below $\varepsilon/4$. The compiler's path variation is
\[
 \widetilde O_d(\kappa^d R_0^{d+1}\varepsilon^{-d}),
\]
where the suppressed factors are powers of $A_\eta=O(\log(1/\varepsilon))$. Its width has the same polynomial bound as before, again up to logarithms.

At widths that fit this compiler, the Gaussian implicit certificate gives the same sample polynomial. At smaller widths $N\ge C+1$, the original comparison function belongs to the learner's tanh-output class, after compiling biases into a constant wire and padding identity hidden layers. The finite-parameter certificate therefore completes the same all-width argument. The two certificate events, comparison concentration, and width-confidence allocation are unchanged.
\end{proof}

There is a necessary distinction between comparison classes. A tanh-output width-$C$ network need not represent the same functions as a raw-readout width-$C$ network. The theorem above compares with networks sharing the learner's final tanh. For comparison with arbitrary raw ReLU teachers, the transformed-interpolation argument still applies after first clipping that teacher, but it only directly proves the guarantee once learner width fits the resulting compiler. The all-width theorem for raw teachers uses the hard saturation of Theorem~\ref{thm:all-width-regular-Gaussian}.
